\section*{Abstract}

Die vorliegende Hausarbeit untersucht die Konzeption und prototypische Implementierung eines Bestellanforderungsprozesses mit integriertem Genehmigungsworkflow auf der SAP Business Technology Platform (BTP). Im Mittelpunkt steht die Frage, wie sich durch die Kombination des SAP Cloud Application Programming Model (CAP), SAP Fiori Elements und SAP Build Process Automation ein durchgängiger, cloudnativer Geschäftsprozess realisieren lässt. Auf Basis einer Anforderungsanalyse wird zunächst ein Architekturkonzept entworfen, das eine klare Schichtentrennung zwischen Datenmodell, Geschäftslogik und Benutzeroberfläche vorsieht. Anschließend wird der Prototyp implementiert, der es Fachanwendern ermöglicht, Bestellanforderungen mit Dokumentenanhängen in einer Fiori-Elements-Oberfläche zu erstellen und zur Genehmigung einzureichen. Der Genehmigungsworkflow wird in SAP Build Process Automation modelliert und über eine REST-Schnittstelle bidirektional mit dem CAP-Service gekoppelt. Der Workflow-Kontext umfasst Positionslisten, Anhänge-Metadaten mit direkten Download-URLs sowie Deep Links zur Fiori-Anwendung. Eine Statushistorie protokolliert den vollständigen Lebenszyklus jeder Bestellanforderung. Die Ergebnisse zeigen, dass die gewählte Architektur eine effiziente Entwicklung ermöglicht und die Plattformdienste der SAP BTP eine nahtlose Integration von Anwendungslogik und Prozessautomatisierung erlauben.

\bigskip
\large{\textbf{Abstract (English)}}

\normalsize
This paper examines the design and prototypical implementation of a purchase requisition process with an integrated approval workflow on the SAP Business Technology Platform (BTP). The central question is how a cloud-native end-to-end business process can be realized by combining the SAP Cloud Application Programming Model (CAP), SAP Fiori Elements, and SAP Build Process Automation. Based on a requirements analysis, an architecture concept is developed that provides a clear separation of concerns between data model, business logic, and user interface. Subsequently, a prototype is implemented that enables business users to create purchase requisitions with document attachments in a Fiori Elements interface and submit them for approval. The approval workflow is modeled in SAP Build Process Automation and bidirectionally coupled with the CAP service via a REST interface. The workflow context includes item lists, attachment metadata with direct download URLs, and deep links to the Fiori application. A status history records the complete lifecycle of each purchase requisition. The results demonstrate that the chosen architecture enables efficient development and that the platform services of SAP BTP allow seamless integration of application logic and process automation.
